% Source: https://artofproblemsolving.com/community/c3h2167243_squaring_manipulation
% Stable identity: aops-topic-2167243-q01
\textbf{(a)} Put $u=x-p$ and $v=y-q$. Then
\[uv=r,\qquad u+v=s-p-q=:t,\]
so $u$ and $v$ are the two roots of $z^2-tz+r=0$ (real roots need $t^2\ge4r$). Hence $(u-v)^2=(u+v)^2-4uv=t^2-4r$ and $u-v=\pm\sqrt{t^2-4r}$. Since $x-y=(u+p)-(v+q)=(u-v)+(p-q)$,
\[\boxed{x-y=p-q\pm\sqrt{(s-p-q)^2-4r}} .\]
Both signs occur when $(s-p-q)^2>4r$: the line $x+y=s$ meets the hyperbola in two points, obtained from each other by swapping the values of $u$ and $v$.

\textbf{(b)} If $p=q$, then $x-y=\pm\sqrt{(s-2p)^2-4r}$, so
\[|x-y|=\sqrt{(s-2p)^2-4r},\]
whichever of the two points $P$ is.

If $p\ne q$ the two values $p-q\pm\sqrt{(s-p-q)^2-4r}$ usually have different absolute values. For example, take $p=1$, $q=0$, $r=2$, $s=4$: then $u+v=3$ and $uv=2$, so $\{u,v\}=\{1,2\}$, and the two points are $P=(2,2)$ and $P=(3,1)$, with $|x-y|=0$ and $|x-y|=2$ respectively.

\begin{asy}
usepackage("amsmath");
size(5.5cm);
defaultpen(fontsize(10pt));
import graph;
xaxis(-2.5,6.5,Arrow(TeXHead)); yaxis(-4.5,6.5,Arrow(TeXHead));
draw(graph(new real(real x){return 2/(x-1);},1.32,6.3),blue+linewidth(0.9));
draw(graph(new real(real x){return 2/(x-1);},-2.4,0.56),blue+linewidth(0.9));
draw((1,-4.5)--(1,6.5),gray+dashed);
draw((-1.5,5.5)--(6,-2),red);
dot((2,2)); dot((3,1));
label("$(2,2)$",(2,2),NE); label("$(3,1)$",(3,1),NE);
label("$x+y=4$",(6,-2),S,red);
label("$(x-1)y=2$",(1.45,4.4),E,blue);
label("$x=1$",(1,-4.5),E,gray);
label("\mbox{}",(min(currentpicture,user=true)+max(currentpicture,user=true))/2,black); // keeps the text after the figure black when LaTeX typesets the labels inline
\end{asy}

So in general the positive difference between the coordinates is not determined by $p$, $q$, $r$ and $s$ alone.
